\clearpage
\SourceLeaf{495}
\pdfbookmark[1]{Nachlass V. Allgemeine Störungen der Pallas durch Jupiter. Erste Rechnung.}{pallas-V}
\fancyhead[L]{\small NACHLASS. STÖRUNGEN DER PALLAS}
\begin{center}
{\Large NACHLASS.\par}
\bigskip
{\large [V. ALLGEMEINE STÖRUNGEN DER PALLAS DURCH JUPITER. ERSTE RECHNUNG.]\par}
\medskip
[BEGONNEN 1811 AUGUST.]\par
\end{center}
\par\Needspace{9\baselineskip}\bigskip
\pdfbookmark[2]{Allgemeine Störungen, Erste Rechnung, [1]}{pallas-V-fragment1}
\begin{center}[1.]\end{center}
\begingroup
\makeatletter\tagsleft@true\let\veqno\@@leqno\makeatother
\setlength{\lineskip}{4pt}\setlength{\lineskiplimit}{4pt}
\newcommand{\VonePallasMarks}{\makebox[.16em][l]{\raisebox{1.35ex}{\rule{.07em}{.07em}}}\kern-.16em\makebox[.16em][l]{\raisebox{-.5ex}{\rule{.14em}{.035em}}}}
\noindent\hspace*{\parindent}[Zur Entwicklung der allgemeinen Theorie der Störungen der Pallas durch Jupiter wurde im Großen und Ganzen, wie bei der Rechnung der speziellen Störungen, die in der »Exposition etc.« (Seiten 449 und 465 dieses Bandes) dargelegte Methode angewandt. Da indessen hier beabsichtigt wird, die allgemeine Theorie darzustellen und nicht nur die numerischen Werte der Störungen für einen bestimmten Zeitraum zu berechnen, müssen die Gleichungen für die Störungen hier in etwas anderer Form aufgestellt werden als dort. Diese sollen darum hier noch einmal zusammengestellt werden, wobei die Bezeichnungsweise dieselbe bleiben soll wie bei den früheren Rechnungen (Seite 473 dieses Bandes).
\par
Die Koordinaten von Pallas und Jupiter \(r,v,r',w'\) lassen sich aus den elliptischen Elementen nach bekannten Formeln berechnen. Zur Berechnung der gegenseitigen Lage der Bahnen dienen ferner die Formeln (vgl. Seite 459 dieses Bandes):
\[
\begin{aligned}
\sin\tfrac12\mathrm{J}.\sin\tfrac12(\Delta+N)&=\sin\tfrac12(i'+i).\sin\tfrac12(\AscendingNode'-\AscendingNode)\\
\sin\tfrac12\mathrm{J}.\cos\tfrac12(\Delta+N)&=\sin\tfrac12(i'-i).\cos\tfrac12(\AscendingNode'-\AscendingNode)\\
\cos\tfrac12\mathrm{J}.\sin\tfrac12(\Delta-N)&=\cos\tfrac12(i'+i).\sin\tfrac12(\AscendingNode'-\AscendingNode)\\
\cos\tfrac12\mathrm{J}.\cos\tfrac12(\Delta-N)&=\cos\tfrac12(i'-i).\cos\tfrac12(\AscendingNode'-\AscendingNode),
\end{aligned}
\tag*{1)}
\]
aus denen sich
\begin{center}
\begin{tabular}{@{}r@{\,}p{.82\linewidth}@{}}
\(\mathrm{J}={}\)&Neigung beider Bahnebenen,\\
\(\AscendingNode+\Delta={}\)&Länge des Knotens der Jupiterbahn auf der \VonePallasMarks Pallasbahn, gezählt auf der letzteren,\\
\(\AscendingNode'+N={}\)&»\quad »\quad »\quad »\quad »\quad »\quad »\quad »\quad »\quad »\quad »\quad ersteren
\end{tabular}
\end{center}
ergeben. Für die Berechnung von \(v'\) und \(\beta'\) hat man weiterhin die Gleichungen
\[
\begin{aligned}
\operatorname{tang}(v'-\AscendingNode-\Delta)&=\cos\mathrm{J}.\operatorname{tang}(w'-\AscendingNode'-N)\\
\operatorname{tang}\beta'&=\operatorname{tang}\mathrm{J}.\sin(v'-\AscendingNode-\Delta)\\
\sin\beta'&=\sin\mathrm{J}.\sin(w'-\AscendingNode'-N),
\end{aligned}
\tag*{2)}
\]
\SourceLeaf{496}von denen eine der beiden letzteren zur Kontrolle dienen kann. Ferner findet sich
\[
\begin{aligned}
\xi&=r'\cos\beta'\cos(v-v')&\quad\rho\rho&=r'r'+rr-2r\xi\\[.4em]
\eta&=r'\cos\beta'\sin(v-v')&\Delta&=\frac{1}{\rho^3}-\frac{1}{r'^3}\\[.4em]
\zeta&=r'\sin\beta'&&\\[.4em]
T&=\frac{\alpha r}{\rho^3}-\alpha\Delta\xi&&\\[.4em]
V&=\alpha r\Delta\eta&\alpha&=\frac{\mu a}{\cos\varphi}\mathpunct{\raisebox{.45ex}{.}}\\[.4em]
W&=\alpha r\Delta\zeta&&
\end{aligned}
\tag*{3)}
\]
Der Wert von \(\alpha\) ist noch mit der Zahl \(206265''\) zu multiplizieren, womit man alle Größen in Bogensekunden erhält.
\par
Nach den Gleichungen 3 auf Seite 474 gelten dann für die Störungen der Neigung und der Knotenlänge die Beziehungen:
\[
\begin{aligned}
\frac{\mathrm{d}i}{n\mathrm{d}t}&=\cos(v-\AscendingNode).W\\[.5em]
\frac{\mathrm{d}\AscendingNode}{n\mathrm{d}t}&=\frac{1}{\sin i}\sin(v-\AscendingNode).W.
\end{aligned}
\tag*{4)}
\]
\par
Für die Störungen des Perihels ergibt eine leichte Transformation der genannten Gleichungen:
\[
\begin{aligned}
\frac{\mathrm{d}\tilde\omega}{n\mathrm{d}t}
&=(1-\cos i)\frac{\mathrm{d}\AscendingNode}{n\mathrm{d}t}
+\frac{p}{e}\cos(v-\tilde\omega).T\\
&\qquad-\frac{1}{e}\{2+e\cos(v-\tilde\omega)\}\sin(v-\tilde\omega).V.
\end{aligned}
\tag*{5)}
\]
\par
Für den Halbparameter \(p\), der an die Stelle der Exzentrizität treten soll, ergibt Gleichung 63 der Exposition unter Berücksichtigung der veränderten Bedeutung von \(V\):
\[
\frac{\mathrm{d}p}{n\mathrm{d}t}=-2a\cos\varphi^2.V,
\]
also
\[
\frac{\mathrm{d}\log p}{n\mathrm{d}t}=-2V.
\tag*{6)}
\]
\par
Des Weiteren ergeben die Gleichungen 3 auf Seite 474
\[
\begin{aligned}
\frac{\mathrm{d}\log\operatorname{hyp} n}{n\mathrm{d}t}
&=3ae\sin(v-\tilde\omega).T+\frac{3aa\cos\varphi^2}{rr}V\\[.7em]
\frac{\mathrm{d}\varepsilon}{n\mathrm{d}t}
&=\operatorname{tang}\tfrac12 i\sin(v-\AscendingNode).W\\
&\quad+\{2r\cos\varphi+p\operatorname{tang}\tfrac12\varphi\cos(v-\tilde\omega)\}T\\
&\quad-\operatorname{tang}\tfrac12\varphi(2+e\cos(v-\tilde\omega))\sin(v-\tilde\omega).V\\[.5em]
&=(1-\cos\varphi)\left(\frac{\mathrm{d}\tilde\omega}{n\mathrm{d}t}
 -(1-\cos i)\frac{\mathrm{d}\AscendingNode}{n\mathrm{d}t}\right)\\
&\quad+(1-\cos i)\frac{\mathrm{d}\AscendingNode}{n\mathrm{d}t}
 +2r\cos\varphi.T.
\end{aligned}
\tag*{7)}
\]
\par
Legt man nun der ganzen Rechnung sowohl für Pallas als auch für Jupiter jeweils ein System elliptischer Elemente zugrunde, so sind die Größen \(T,V,W\) ebenso wie die Ableitungen der Elemente offenbar periodische Funktionen der mittleren Anomalien beider Planeten, also der Größen \(M\) und \(M'\), und lassen sich durch trigonometrische Reihen in Vielfachen dieser beiden Argumente darstellen. Man kann darum hier die in der Theoria interpolationis methodo nova tractata (Band III, Seite 265) angestellten Untersuchungen verwenden und die Koeffizienten dieser Reihen bestimmen, indem man die numerischen Werte von \(T,V,W\) für Wertreihen von \(M\) und \(M'\) berechnet, die jeweils eine ganze Periode umfassen.
\par
\SourceLeaf{497}Zu diesem Zweck sollen die in der genannten Abhandlung abgeleiteten Formeln in den folgenden Artikeln so umgeformt werden, dass sie für die Lösung der vorliegenden Aufgabe möglichst einfach sind.]
\par
\endgroup
